% GENERATED, edit the French unit or the translation table
% GENERATED from content/data/enrichment.json by mxbuild.enrichment; do not edit.
\mxunit{enrichment-pigeonhole-notes}{The pigeonhole principle}
\mxsubtitle{High school · olympiad thinking \textperiodcentered{} Lesson and method \textperiodcentered{} 02}

\begin{prerequisites}
Euclidean division and elementary counting.
\end{prerequisites}

\begin{keypoint}[Ideas and method]
If $n + 1$ objects are placed in n boxes, at least one box contains two objects. Otherwise, each box would contain at most one, giving at most n objects.

More generally, if N objects occupy k boxes, one box contains at least $\lceil \frac{N}{k}\rceil $. The ceiling $\lceil x\rceil $ is the smallest integer greater than or equal to x.

Choosing the boxes is the key: remainders, birth months, small grid squares or classes of integers. Explain why every object belongs to one box.
\end{keypoint}

\begin{example}
Among 13 people, two were born in the same month: the 12 months are the boxes. This does not guarantee two birthdays on the same day.
\end{example}

\begin{activity}[Explore]
Choose six distinct integers from 1 to 10. Group them into pairs summing to 11 and find a complete pair.
\end{activity}

